\section{Experimental Design and Evaluation}
\label{sec:experiments}

We evaluate LOKA in the MuJoCo physics simulation environment on an underactuated bipedal walker driven by MuJoCo MPC (MJPC), with planned extensions to additional humanoid embodiments.

\subsection{Evaluation Scenarios}
\label{sec:scenarios}

To validate robust reformulation across physical and semantic shifts, the system is subjected to four validation tasks:
\begin{itemize}
\item \textbf{Mass Shift:} An unmodeled mass of \(x\)\,kg (TODO) is attached asymmetrically mid-stride, forcing the system to detect effort imbalance and re-stabilize its center-of-mass target.
\item \textbf{Friction Shift:} The global surface friction coefficient is stepped down from a nominal value of $1.0$ to a reduced value (TODO) to evaluate real-time environmental adaptation.
\item \textbf{Hardware Failure:} A hard fault is injected into a critical actuator (e.g., right hip), setting physical torque capacity to zero to validate the virtual-amputation pipeline.
\item \textbf{Semantic / Operator Objectives:} The robot receives overarching language instructions (e.g., crouch and walk, crawl, hold upright) and/or must navigate a specialized corridor featuring a low overhanging ceiling followed by elevated rubble.
\end{itemize}

\subsection{Quantitative Performance Metrics}
\label{sec:metrics}

The framework is evaluated using the following metrics:
\begin{itemize}
\item \textbf{Task Success Rate (TSR):} Fraction of trials where the robot reaches the terminal criterion without falling.
\item \textbf{Time to Stabilization ($T_s$):} Duration between anomaly onset and restoration of a steady-state gait (or satisfaction of updated mission criteria).
\item \textbf{Effort Deviation Integral ($E_{dev}$):} Accumulated energy metric in Equation~\ref{eq:effort}, quantifying adaptation efficiency (alternative integrals over tracking residuals are also under consideration).
\item \textbf{Outer-loop turns:} Number of orchestrator interventions until success or timeout.
\item \textbf{Semantic Accuracy Percentage (SAP):} An automated judge (independent frontier model) comparing parsed scratchpad mutations against hidden fault metadata.
\item \textbf{Interface reliability:} Rates of invalid YAML and unresolved MJCF names.
\end{itemize}

\subsection{Comparative Baselines}
\label{sec:baselines}

\textbf{TODO:} Flesh out precise baseline implementations adapted to the dynamic bipedal domain (e.g., vanilla fixed MJPC, Domain Randomization without online mutation, ablations without model mutations or multi-turn history, and a code-generation / Language-to-Rewards variant where latency permits).

\begin{table}[t]
  \centering
  \caption{Placeholder results (fault recovery \%, operator success \%, median outer-loop turns).}
  \label{tab:results}
  \setlength{\tabcolsep}{3.5pt}
  \begin{tabular}{@{}lccc@{}}
    \toprule
    Method & Fault (\%) & Op.\ (\%) & Turns \\
    \midrule
    Fixed MJPC & -- & -- & -- \\
    Ablation (no mut.) & -- & -- & -- \\
    LOKA (full) & -- & -- & -- \\
    \bottomrule
  \end{tabular}
\end{table}

% \begin{figure}[t]
%   \centering
%   \includegraphics[width=\linewidth]{recovery_curves.pdf}
%   \caption{Tracking error vs.\ time after fault injection for LOKA and baselines.}
%   \label{fig:recovery}
% \end{figure}
